%lezione 6 del 20 marzo 2026
% Qui ci sarebbero degli esercizi sulle matrici densità che non ho riportato
% Ometto anche l'esempio sullo spin 1/2 perché ne ho le palle piene
\subsubsection{Esempio: la polarizzazione della luce}
Classicamente, la polarizzazione della luce è definita dal vettore campo elettrico
$$\mathbf{E}(z,t) = 
\begin{pmatrix}
    E_1\cos(kz-\omega t+\alpha_1)\\
    E_2\sin(kz-\omega t +\alpha_2)
\end{pmatrix} = \sqrt{E_1^2+E_2^2}\Re\left[
\begin{pmatrix}
    a_1\\
    a_2
\end{pmatrix}
e^{i(kz-\omega t)}
\right] $$
con $a_i = \frac{E_i}{\sqrt{E_1^2+E_2^2}}e^{i\alpha_i}$ tali che $|a_1|^2+|a_2|^2 = 1$.

Se $\frac{a_2}{a_1}\in\mathbb{R}$, la polarizzazione è lineare in direzione
$\theta=\atan\left(\frac{a_2}{a_1}\right)$; se invece $\frac{a_2}{a_1} = \pm i$, la polarizzazione è
circolare.

In teoria quantistica dei campi, i fotoni sono particelle con $s=1$ e il loro spazio di Hilbert sarebbe
a $2s+1=3$ dimensioni, se non fosse che, essendo \textit{massless}, essi possiedono simmetria solo attorno
all'asse del proprio momento, riducendo così la dimensione a 2.

Lo stato di un fotone può dunque essere descritto come un vettore $\ket{\psi}\in\mathcal{H}=\mathbb{C}^2$.
Una possibile base di questo spazio è $\left\{\ket{x},\ket{y}\right\}$, che rappresentano
(rispettivamente) la polarizzazione lungo $x$ e $y$ e un qualunque stato puro del fotone può essere
rappresentato in questa base come
$$\ket{\psi}=a_1\ket{x}+a_2\ket{y}\text{, con }a_1,a_2\in\mathbb{C}\text{ e }|a_1|^2+|a_2|^2=1$$
La sua rappresentazione come matrice densità è invece
$$\rho_\psi = \ket{\psi}\bra{\psi} = 
\begin{pmatrix}
    |a_1|^2 & a_1a_2^*\\
    a_2a_1^*&|a_2|^2
\end{pmatrix}
$$
Come prima, la polarizzazione è lineare di angolo $\theta=\atan\left(a_2/a_1\right)$ se
$a_2/a_1\in\mathbb{R}$ ed è in generale ellittica se $a_2/a_1 \notin\mathbb{R}$.

Questi stati puri hanno delle naturali rappresentazioni come matrici densità:
$$\theta = \pi/4 \rightarrow\rho_{\pi/4} = \frac{1}{2}
\begin{pmatrix}
    1&1\\
    1&1
\end{pmatrix}
$$
$$
\text{Circolare}\circlearrowleft \rightarrow \rho_\circlearrowleft = \frac{1}{2}
\begin{pmatrix}
    1&-i\\
    i&1
\end{pmatrix}
$$
Per rappresentare la luce \textit{non} polarizzata, si utilizzano stati misti.

\textbf{Esempio} Un fotone che ha la stessa probabilità di essere polarizzato lungo $x$ e lungo $y$
è descritto da
$$\rho=\frac{1}{2}\left(\ket{x}\bra{x}+\ket{y}\bra{y}\right)$$
che, nella base $\left\{\ket{x},\ket{y}\right\}$ non è altro che $\rho=\frac{1}{2}\mathbb{I}$.

Si consideri inoltre lo stato puro 
$$\ket{\psi}=\frac{1}{\sqrt{2}}\left(\ket{x}+\ket{y}\right)$$
come si era visto, la probabilità di trovare la luce polarizzata lungo $x$ è la medesima:
$$\mathcal{P}_\rho(x=1) = Tr(\rho P_{x}) = Tr\left(\frac{1}{2}\mathbb{I}\ket{x}\bra{x}\right) =
\frac{1}{2}$$
$$\mathcal{P}_\psi(x=1) = Tr(\ket{\psi}\bra{\psi}\ket{x}\bra{x}) = \bra{\psi}\ket{x}\bra{x}\ket{\psi} = 
\frac{1}{2}$$
Ma se invece si calcola la probabilità di trovare la luce polarizatta a $\theta=\frac{\pi}{4}$ si trova
$$\mathcal{P}_\psi(\theta=\pi/4) = Tr\left(\ket{\psi}\bra{\psi}\ket{\frac{\pi}{4}}\bra{\frac{\pi}{4}}\right) = 1
\quad\text{(poiché $\ket{\psi}=\ket{\frac{\pi}{4}}$)}
$$
$$\mathcal{P}_\rho(\theta=\pi/4) = Tr(\rho P_{\frac{\pi}{4}}) = \frac{1}{2}Tr(P_{\frac{\pi}{4}}) =
\frac{1}{2}$$
questo vale $\forall \theta$, in quanto la traccia di un proiettore è sempre 1.

